\subsubsection*{Statement B.2m (Proved): the kernel bound with arbitrary alteration constants}
\begin{proposition}
\label{B:B2-alteration-constants}
Let $X_0$ be a proper geometrically normal geometrically integral
$k$-scheme of characteristic $p>0$.  We may choose a finite field
extension $l/k$, a smooth projective geometrically integral
$l$-scheme $Y$, and a proper generically finite $k$-morphism
$Y\to X_0$.  Put $d=[k(Y):k(X_0)]$ and
\[
 P=\operatorname{Pic}^{\tau}_{Y/l},\qquad
 \mathcal R=\operatorname{Res}_{l/k}(P),
 \quad\mathcal R(T)=P(T\times_k l).
\]
For each proper geometrically normal modification $X\to X_0$,
there is a compatible homomorphism of fppf sheaves
\[
 h_X:\operatorname{Pic}^{\tau}_{X/k}\longrightarrow\mathcal R
\]
whose kernel is killed by $d$.  There is an integer $g$ depending
only on $k(X_0)/k$ such that all these kernels are finite of rank
at most $d^{2g}$.  Thus a fixed comparison sheaf and a uniform
kernel bound are available without assuming that an alteration
is geometrically integral over the original field.
\end{proposition}
\begin{proof}
For existence, apply the alteration theorem over an algebraic
closure of $k$, and descend the smooth projective geometrically
integral source and its proper generically finite morphism to a
finite extension $l/k$, enlarging $l$ if necessary.  Smoothness,
geometric integrality and projectivity descend, and the finite
generic degree can be descended on an open where the map is finite
locally free.  This uses the projective alteration supplied by
\cite[Corollary 5.15, with the projectivity assertion in
Theorem 5.9 and Proposition 5.11]{deJong1997} and finite-presentation
descent \cite[Tag 01ZM]{StacksB}.  No separability assertion about
$l/k$ is required.

Let $Y_X$ be the normalization of the main component of
$Y\times_{X_0}X$.  The projection $Y_X\to Y$ is a proper
birational morphism of $l$-schemes.  Its source is normal and
geometrically integral over $l$: its function field is $k(Y)$,
which is a regular extension of $l$.  Statement B.1g over $l$
therefore identifies
\[
 \operatorname{Pic}^{\tau}_{Y_X/l}=P.
\]
Pullback from $X_l$ to $Y_X$, followed by this identification,
defines $h_X$.  Pullback preserves the $\tau$-part, so its target
is the indicated Weil restriction.  Formation of Weil restriction
commutes with base change, and $\mathcal R$ is an algebraic space
\cite[Lemma 97.11.2 and Proposition 97.11.5, Tag 05Y8]{StacksB}.
It is separated: restriction of scalars along a finite locally
free map preserves a closed immersion, as is checked by the
finitely many coefficient equations of its ideal, and we apply
this to the diagonal of $P$.  In particular the kernels under
discussion are closed subgroup schemes of
$\operatorname{Pic}^{\tau}_{X/k}$.

We verify that passing from relative Picard classes over $k$ to
relative classes over $l$ has not enlarged the kernel.  Work fppf
locally on a $k$-scheme $T$ so that a section of
$\operatorname{Pic}^{\tau}_{X/k}$ is represented by a line bundle
$\mathcal L$ on $X_T$.  Write $T_l=T\times_k l$.
If $h_X(\mathcal L)=0$, then the pullback of $\mathcal L$ to
$(Y_X)_T$ has zero class in the relative Picard sheaf for
$(Y_X)_T\to T_l$.  Since
$\mathcal O_{T_l}=q_*\mathcal O_{(Y_X)_T}$ universally, the
low-degree Leray sequence for $\mathbf G_m$ says that this line
bundle is pulled back from a line bundle $\mathcal M$ on $T_l$.

Such an $\mathcal M$ becomes trivial after an fppf cover of $T$,
even if $l/k$ is inseparable and $T$ is nonreduced.  Indeed let
$\pi:T_l\to T$.  The functor of generators of $\mathcal M$ as
an $\mathcal O_{T_l}$-module is an open subscheme of the vector
bundle associated with the finite locally free
$\mathcal O_T$-module $\pi_*\mathcal M$.
The open condition is that the determinant of
$\mathcal O_{T_l}\to\mathcal M$, $1\mapsto m$, is invertible.
It is smooth over $T$ and has a point over every geometric
fibre, since an invertible module over a finite product of Artin
local rings is free of rank one.  Hence it is a smooth surjective,
and therefore fppf, cover of $T$ trivializing $\mathcal M$.
Consequently $\ker h_X$ is exactly the kernel of the ordinary
$k$-relative Picard pullback along $Y_X\to X$.

The latter proper dominant map has generic degree $d$ and normal
source and target.  Statement B.2j therefore kills its entire
fppf Picard kernel by $d$.  For the rank bound, extend to an
algebraic closure and then to the fixed generic-curve field used
in B.2d.  If $X_0$ is not projective, take the normalization of
the graph between $X$ and a fixed projective normal model before
restricting to that curve.  Pullback from $X$ to the graph is a
closed immersion on $\operatorname{Pic}^{\tau}$, and B.2d embeds
the latter into one fixed $g$-dimensional Jacobian $J$.
The kernel is therefore a closed subgroup of $J[d]$ and has
rank at most $d^{2g}$.  Finiteness and this rank bound descend
through the field extensions.
\end{proof}

\begin{remark}
We have not identified the $\tau$-part of the Picard scheme of
the possibly geometrically nonreduced $k$-scheme $Y$ with a Weil
restriction.  The kernel calculation only uses the relative
Picard scheme over $l$, the explicit fppf trivialization of base
line bundles, and the norm argument of B.2j.  The missing uniform
bound for the images $h_X$ is unchanged.
\end{remark}
